\chapter{С планковского пола вверх}
\label{ch:floors-preface}

В \cref{sec:planck-floor} мы спустились до планковского мира~$\Mlayer$:
счётные узлы~$\Lambda$, такт~$\htick$, ячейка~$\hv=\PlanckCell$, закон~$g$.
Дальше текст идёт не вниз по единицам измерения, а \emph{вверх по радиусу} ---
на том же носителе, без новых аксиом.

\paragraph{Три этажа.}
Перед инфракрасной hop-лестницей ($r_e$, комптоновская длина, боровский радиус;
\cref{ch:si-sm}) на~$\Mlayer$ выделяются три этажа составного узла:

\begin{enumerate}[label=\textbf{этаж~\arabic*.},leftmargin=*]
\item \textbf{Планковский пол} ($N\sim 1$): один планкон в одной ячейке~$\hv$.
  Внутренний спектр --- фазовое кольцо $\mathbb{Z}_{512}$ и спинор $SU(2)$ на ядре
  (\cref{ch:floor0}).

\item \textbf{Чуть выше} ($N\sim N_{12}\ldots N_{12}^{5}$): облако~$\rho_\Theta$
  вокруг ядра, $\varepsilon$-звезда соседей, минимальный ореол $R_{\mathrm{dress}}=\hl$,
  полоса пре-резонансов на степенях каузальной звезды
  (\cref{ch:floor1}).

\item \textbf{Составной узел} ($N\sim 10^{15}\cdot\hl$): конфайнящие фокусы,
  кварковая схема, квантовые числа цвета и поколения как следствия топологии,
  не как вход (\cref{ch:floor2}).
\end{enumerate}

\paragraph{Что не смешивать.}
Этаж~1 --- не комптоновский радиус и не боровская орбиталь
(те --- $k\approx 8\ldots 9$ на другой оси, \cref{prop:ladder-scale-steps}).
Этаж~0 --- не «электрон целиком'': лабораторный электрон --- readout на~$T$
после всех этажей и IR-моста.
Этаж~2 --- ещё на~$\Mlayer$; гладкая квантовая химия --- дальше по лестнице~$k\geq 10$.

\paragraph{Порядок чтения.}
Главы~\ref{ch:floor0}--\ref{ch:floor2} можно читать подряд как одну прогулку вверх
от одной занятой ячейки к составной материи.
Общая механика планкона и жидкости остаётся в \cref{ch:matter};
константы и полная лестница до метра --- в \cref{ch:si-sm,ch:alpha}.
